% =============================================================================
%  Carta de Apresentação, Avelino Almeida, 2026
%  Compilar no Overleaf com XeLaTeX. Para compilar a carta em vez do CV,
%  definir este ficheiro como principal (Menu > Main document > carta.tex).
%  Por candidatura: preencher dados/candidatura.tex e o parágrafo
%  "Porquê esta empresa" abaixo.
% =============================================================================
\documentclass[11pt]{article}
\usepackage[a4paper,left=2cm,right=2cm,top=1.4cm,bottom=1.6cm]{geometry}
\usepackage{estilo/cvestilo}
% =============================================================================
%  Dados pessoais: usados no cabeçalho do CV e da carta.
% =============================================================================
\newcommand{\CVnome}{Avelino}
\newcommand{\CVapelido}{Almeida}
\newcommand{\CVtitulo}{Consultor DBA \textperiodcentered{} SQL Server, BI e ETL}
\newcommand{\CVlocal}{Penafiel, Portugal}
\newcommand{\CVemail}{avlalmeida@gmail.com}
\newcommand{\CVtelefone}{+351 916 472 603}
\newcommand{\CVlinkedin}{linkedin.com/in/avelinoalmeida}
\newcommand{\CVsite}{avelinoalmeida.github.io/home}

% =============================================================================
%  Dados da candidatura: o único ficheiro a mudar em cada candidatura,
%  juntamente com o parágrafo "Porquê esta empresa" em carta.tex.
%  Os campos por preencher aparecem a vermelho no PDF.
% =============================================================================
\newcommand{\empresa}{\pf{Nome da empresa}}
\newcommand{\destinatario}{\pf{Equipa de Recrutamento ou nome da pessoa}}
\newcommand{\moradaempresa}{\pf{Morada da empresa}}
\newcommand{\funcao}{\pf{Nome da função}}
\newcommand{\referencia}{\pf{Referência da oferta}}
\newcommand{\saudacao}{Exmos. Senhores,}   % ou: Caro(a) Dr.(a) Apelido,


\hypersetup{pdftitle={Avelino Almeida - Carta de Apresentação},pdfauthor={Avelino Almeida}}
\setlength{\parskip}{7pt}
\fancyfoot[C]{\footnotesize\color{cinzaclaro}Avelino Almeida \textperiodcentered{} Carta de Apresentação}

\color{texto}
\begin{document}
\raggedright
\cabecalho
\vspace{6pt}

% --- Destinatário, data e assunto --------------------------------------------
{\robotobold \destinatario}\\
{\color{cinza}\empresa\\ \moradaempresa}

\vspace{2pt}
{\color{cinza}Penafiel, \datapt}

{\robotobold Assunto:} Candidatura a \funcao{} (ref.\ \referencia)

\saudacao

% --- Corpo -------------------------------------------------------------------
\secao{Sobre mim}
Sou Consultor DBA na eSPap, onde administro ambientes SQL Server on-premises e em Azure, integrado numa equipa de três DBA responsável por um parque de grande escala da Administração Pública. O meu dia a dia passa pela manutenção preventiva, pela monitorização e análise de desempenho, pela gestão de cópias de segurança e recuperação de dados, e pela administração dos servidores de SSIS, SSAS e SSRS.

Antes disso, na ULS Tâmega e Sousa, desenvolvi um protótipo de Business Intelligence para o projeto-piloto de saúde mental CRI-SM, do data warehouse em SQL Server ao modelo dimensional e aos painéis em Power BI. Em paralelo, sou formador certificado há cinco anos e frequento o Mestrado em Engenharia Informática no Politécnico do Porto.

\secao{Porquê esta empresa}
\pf{Duas ou três frases específicas para esta candidatura: o que conhece da empresa, o projeto, o setor ou a tecnologia referidos na oferta, e porque encaixam no seu percurso. Exemplo de estrutura: «Acompanho o trabalho da empresa em X. A oferta refere Y, que corresponde ao que faço hoje com Z.»}

\secao{Porquê eu}
Trago três contributos. Experiência operacional em SQL Server de produção, num contexto em que a disponibilidade e a recuperação dos dados são críticas. Capacidade de levar os dados até à decisão, da modelação dimensional aos dashboards em Power BI. E o hábito de explicar temas técnicos a públicos diferentes, consolidado em mais de 1000 horas de formação.

A isto junto 17 anos na indústria, sete deles a coordenar uma equipa, onde aprendi a gerir prazos, prioridades e pessoas.

Agradeço a atenção dispensada e fico disponível para uma conversa, presencial ou remota, quando for conveniente.

Com os melhores cumprimentos,

\vspace{14pt}
{\robotobold \CVnome{} \CVapelido}

\vfill
{\small\color{cinza}{\robotobold Anexo:} Curriculum Vitae}

\end{document}
